% Generated draft: requires mathematical and source review.
\documentclass[11pt]{article}
\usepackage[margin=25mm]{geometry}
\usepackage{fontspec}
\setmainfont{Latin Modern Roman}
\usepackage{amsmath,amssymb,mathtools,unicode-math}
\setmathfont{Latin Modern Math}
\usepackage{enumitem,needspace,xurl,hyperref}
\setlength{\parindent}{0pt}
\setlength{\parskip}{5pt}
\setlength{\emergencystretch}{3em}
\newcommand{\N}{\mathbb{N}}
\newcommand{\Z}{\mathbb{Z}}
\newcommand{\Q}{\mathbb{Q}}
\newcommand{\R}{\mathbb{R}}
\newcommand{\C}{\mathbb{C}}
\newcommand{\F}{\mathbb{F}}
\newcommand{\A}{\mathbb{A}}
\newcommand{\PP}{\mathbb{P}}
\newcommand{\E}{\mathbb{E}}
\newcommand{\eps}{\varepsilon}
\newcommand{\dd}{\,\mathrm{d}}
\begin{document}
\section*{UnsolvedMath ID: 11100004}
\textbf{Source problem code:} catalogue entry
\subsection{Definitions and mathematical statement}
Hovey's major-problem item 4 asks for a way to do index theory related to elliptic cohomology. The source does not select an elliptic cohomology theory, analytic operator class, manifold/family category or theorem to prove. In a usual index-theoretic formulation one would need an analytic construction with values in a chosen theory \(E\), appropriate orientation/pushforward data and a comparison with a topological index; the record supplies none of these as fixed requirements. An even-periodic theory associated with an elliptic curve, periodic topological modular forms \(\mathrm{TMF}\), and connective \(\mathrm{tmf}\) are not interchangeable targets by notation alone.

The source proposes possible analytic connections with Richard Melrose's manifolds-with-corners index theory and Ezra Getzler's work on certain infinite-dimensional manifolds. It says these \emph{might} be relevant and expresses limited knowledge; it names Haynes Miller and Mike Hopkins as people who have considered the algebraic-topology side. These are suggestions, not a claimed construction or an index theorem.

The supplied partial-progress summary describes the Witten genus and a \(\sigma\)-orientation as elliptic-index progress, while saying a full \(\mathrm{TMF}\) index theory analogous to Atiyah--Singer is not supplied. Several of its more specific assertions need source correction: the usual string-orientation framework involves a string structure, a refinement of a spin structure trivializing the appropriate \(p_1/2\) class, and a Thom-spectrum map with \(M\operatorname{String}\) as domain. The summary instead writes \(M\operatorname{Spin}\to\mathrm{tmf}\) and attributes to ``Stolz's theorem (1996)'' a vanishing-Witten-genus/orientation implication without hypotheses or a reference. Those statements are not adopted as established facts or as substitutes for the intended orientation. Its phrases ``p-index'' and ``multiplicative/Gray dynamism'' are undefined.

The summary additionally names Ando--Hopkins--Strickland--Rezk maps, Hopkins--Mahowald, Landweber exactness/reflection and \(e\)-invariant connections, then ends with a truncated Bunke remark. No completed analytic construction, exact comparison formula or reliable complete citation can be recovered from that list. The broad elliptic-index program is therefore retained with its supplied uncertainty, rather than declared solved by a chosen genus alone.
\subsection{Short English statement}
Develop index theory linked to elliptic cohomology. The source offers tentative analytic directions rather than a defined operator/index theorem. The supplied TMF/Witten-genus progress summary is retained with its missing references, inconsistent orientation domain and truncation explicitly flagged.
\subsection{Sources}
\begin{itemize}
\item Local UnsolvedMath registry, permanent ID \(11100004\), exact source code \texttt{catalogue entry}, ``Major problems, item 4,'' attributed to Mark Hovey; no proposed year supplied.
\item \emph{Hovey's Algebraic Topology Problem List}, surviving mirror: \url{https://www-users.cse.umn.edu/~tlawson/hovey/}. The background describes 65 top-level items across nine category pages. HTML-ordered-list extraction, accessed 2026-07-29.
\item Source suggestions: Ezra Getzler, Richard Melrose, Haynes Miller and Mike Hopkins, with no complete work references. Their possible relevance is explicitly tentative.
\item Supplied summary names Witten genus, \(\sigma\)-orientation, Ando--Hopkins--Strickland--Rezk, Hopkins--Mahowald, ``Stolz's theorem'' (1996), Landweber and \(e\)-invariant ideas, Atiyah--Singer as a comparison, and a truncated mention of Bunke. No full citations or precise hypotheses are supplied; no missing reference is invented.
\item Registry status partially\_solved and classification PARTIAL-PROGRESS report elliptic-genus/orientation progress without a complete general index theory. Source evidence NEEDS\_REVIEW presents an open problem and requires a current primary-literature audit. Specific summary formulas and theorem attributions remain unverified as discussed.
\item Dataset attribution: Ulam AI / UnsolvedMath Contributors, CC BY 4.0. Source rights NEEDS\_REVIEW: publicly accessible author/publisher copy, redistribution terms requiring review.
\end{itemize}
\end{document}
